% Purpose: Demonstrate bounded notation and a stepwise quadratic derivation.
% Status: Draft template; human mathematical and layout review pending.
% Source-of-truth: .claude/agents/latex-docxer.md (extracted algebraic example).
% Origin: Symbolic algebra only; no project implementation or test fixture.
\documentclass{article}
\usepackage[a4paper,margin=2cm]{geometry}
\usepackage{amsmath,amssymb}

\begin{document}
\section*{Illustrative quadratic derivation}

Let $a,b,c,x\in\mathbb{R}$, with $a\ne0$. Define the discriminant
$\Delta=b^2-4ac$. The real-root derivation below assumes $\Delta\ge0$;
when $\Delta<0$, the equation has no real roots.

\begin{equation}
  \begin{aligned}
    ax^2+bx+c &= 0
      && \text{Given equation} \\
    x^2+\frac{b}{a}x &= -\frac{c}{a}
      && \text{Divide by }a \\
    x^2+2\frac{b}{2a}x+\left(\frac{b}{2a}\right)^2
      &= -\frac{c}{a}+\left(\frac{b}{2a}\right)^2
      && \text{Complete the square} \\
    \left(x+\frac{b}{2a}\right)^2
      &= \frac{\Delta}{4a^2}
      && \text{Collect terms} \\
    x+\frac{b}{2a} &= \pm\frac{\sqrt{\Delta}}{2a}
      && \text{Absorb the sign of }a\text{ into }\pm \\
    x &= \frac{-b\pm\sqrt{\Delta}}{2a}
      && \text{Isolate }x
  \end{aligned}
  \label{eq:quadratic-formula-derivation}
\end{equation}

The two signs in Equation~\ref{eq:quadratic-formula-derivation} coincide
when $\Delta=0$. This example demonstrates algebraic exposition; it does
not assert alignment with a repository algorithm or automated test.
\end{document}
